%----------------------------------------------------------------------------------------
% RATIONALE FUNCTIES EN ASYMPTOTEN. Alle opgaven zijn eigen werk.
%
% Opbouw: Start, leerstof (theorie, werkwijze of uitleg, voorbeeld, korte oefening), en op een
% nieuwe pagina de Opdrachten (door elkaar, stijgende moeilijkheid), Koppel en Test jezelf.
% Labels van dit hoofdstuk beginnen met h2-.
%----------------------------------------------------------------------------------------

\mychapter{2}{Rationale functies}

\lkr{Leerplan B+S': LPD 18 en K1 (rationale functies), het deel over domein, nulwaarden, perforaties, asymptoten en het schetsen van de grafiek. Limieten komen in de cursus Afgeleiden; hier worden asymptoten informeel ingevoerd (\textit{de functiewaarden worden onbeperkt groot}, \textit{het verschil met de rechte wordt willekeurig klein}) en bepaald met de euclidische deling uit hoofdstuk 1.}

Hoe duur is één stuk als je er $q$ maakt, met $500$ euro vaste kosten en $2$ euro per stuk? Dat is $\frac{500 + 2q}{q}$: een quotiënt van twee veeltermen. Zulke functies heten rationaal. Hun grafieken hebben vaak lijnen waar ze steeds dichter bij komen zonder ze te raken: asymptoten.

\begin{oefeningen}[Start]
\oef{h2-start-1}

\oef{h2-verken}
\end{oefeningen}

%========================================================================================
\section{Domein, nulwaarden en perforaties}

\begin{theorie}
Een \textbf{\textcolor{accent}{rationale functie}} is een quotiënt van twee veeltermfuncties: $\formule{f(x) = \frac{A(x)}{B(x)}}$. Het domein bestaat uit alle reële getallen behalve de nulwaarden van de noemer.

Hebben teller en noemer een gemeenschappelijke factor $x - a$, dan mag je die schrappen voor $x \neq a$. De grafiek heeft dan een \textbf{\textcolor{accent}{perforatie}} in $a$: een gat in de grafiek van de vereenvoudigde functie.
\end{theorie}

\voorbeeld{}{Bepaal domein, nulwaarden en perforaties van $f(x) = \dfrac{x^2 - 4}{x^2 - x - 2}$.\label{vb:h2-perforatie}}{%
\uitleg{Ontbind teller en noemer. Het domein lees je af vóór je vereenvoudigt; de nulwaarden en de perforatie erna.}
\begin{stappenbreed}
f(x) = \frac{(x - 2)(x + 2)}{(x - 2)(x + 1)} & Ontbinden. \\
\dom f = \R \setminus \{-1, 2\} & Nulwaarden van de noemer. \\
f(x) = \frac{x + 2}{x + 1} \text{ voor } x \neq 2 & Factor $x - 2$ schrappen. \\
\text{nulwaarde } -2,\ \text{perforatie in } \bigl(2, \tfrac43\bigr) & $\frac{2 + 2}{2 + 1} = \frac43$.
\end{stappenbreed}}

\begin{oefeningen}
\needspace{6\baselineskip}
\oef{h2-domein}
\end{oefeningen}

%========================================================================================
\section{Asymptoten}

\naast{%
Bij $f(x) = \frac{1}{x}$ worden de functiewaarden onbeperkt groot als $x$ dicht bij $0$ komt: $f(0{,}001) = 1000$. De $y$-as is een \textbf{verticale asymptoot}. Voor grote $|x|$ komen de functiewaarden dicht bij $0$: de $x$-as is een \textbf{horizontale asymptoot}.

Elke rationale functie gedraagt zich voor grote $|x|$ zoals het quotiënt $Q(x)$ van de euclidische deling: $f(x) = Q(x) + \frac{R(x)}{B(x)}$, en omdat de rest een lagere graad heeft dan de noemer, wordt $\frac{R(x)}{B(x)}$ willekeurig klein.}{%
\begin{tikzpicture}[x=.45cm,y=.45cm,font=\scriptsize]
  \rooster{-4}{1}{4}{-4}{1}{4}
  \kromme{.25}{4}{1/\x} \kromme{-4}{-.25}{1/\x}
  \node[accent,above right] at (.5,2) {$\frac1x$};
\end{tikzpicture}}

\begin{theorie}
Voor een rationale functie $f(x) = \frac{A(x)}{B(x)}$, na vereenvoudigen:
\begin{itemize}
\item een nulwaarde $a$ van de noemer geeft een \textbf{\textcolor{accent}{verticale asymptoot}} $\formule{x = a}$;
\item graad $A <$ graad $B$: \textbf{\textcolor{accent}{horizontale asymptoot}} $\formule{y = 0}$;
\item graad $A =$ graad $B$: horizontale asymptoot $\formule{y = \frac{a_n}{b_n}}$, het quotiënt van de hoogstegraadscoëfficiënten;
\item graad $A =$ graad $B + 1$: \textbf{\textcolor{accent}{schuine asymptoot}} $\formule{y = Q(x)}$, het quotiënt van de deling $A(x) : B(x)$;
\item graad $A >$ graad $B + 1$: geen horizontale of schuine asymptoot.
\end{itemize}
\end{theorie}

\voorbeeld{}{Bepaal de asymptoten van $f(x) = \dfrac{2x^2 + 1}{x^2 - 1}$.\label{vb:h2-ha}}{%
\uitleg{Teller en noemer hebben geen gemeenschappelijke factor (de teller is nooit nul). Teller en noemer hebben dezelfde graad.}
\begin{stappenbreed}
x^2 - 1 = 0 \iff x = \pm 1 & Verticale asymptoten $x = -1$ en $x = 1$. \\
y = \frac{2}{1} = 2 & Horizontale asymptoot. \\
f(x) = 2 + \frac{3}{x^2 - 1} & Controle met de deling.
\end{stappenbreed}}

\begin{oefeningen}
\needspace{6\baselineskip}
\oef{h2-ha}
\end{oefeningen}

\voorbeeld{}{Bepaal de asymptoten van $f(x) = \dfrac{x^2 + x - 1}{x - 1}$.\label{vb:h2-sa}}{%
\uitleg{De graad van de teller is één meer dan die van de noemer: deel met Horner.}
\begin{stappenbreed}
f(x) = x + 2 + \frac{1}{x - 1} & Horner met $a = 1$: $1, 2, 1$. \\
x = 1 & Verticale asymptoot. \\
y = x + 2 & Schuine asymptoot: het verschil $\frac{1}{x - 1}$ wordt klein.
\end{stappenbreed}}

\begin{oefeningen}
\needspace{6\baselineskip}
\oef{h2-sa}

\oef{h2-asymptoten}

\needspace{6\baselineskip}
\oef{h2-va}
\end{oefeningen}

%========================================================================================
\section{De grafiek schetsen}

\voorbeeld{}{Schets de grafiek van $f(x) = \dfrac{x - 1}{x^2 - 4}$.\label{vb:h2-schets}}{%
\naast{%
\uitleg{Verzamel alle informatie en zet ze in een tekentabel; die zegt waar de grafiek boven of onder de $x$-as ligt, ook naast de asymptoten.}
\begin{stappenvrij}{4cm}{3.3cm}
\dom f = \R \setminus \{-2, 2\} & Geen perforaties. \\
\text{nulwaarde } 1 & \\
x = -2,\ x = 2 & Verticale asymptoten. \\
y = 0 & Graad teller $<$ graad noemer. \\
f(0) = \tfrac14 & Snijpunt met de $y$-as.
\end{stappenvrij}
\figuren{\tekentabel{cccccccc}{x & & -2 & & 1 & & 2 & }{f(x) & - & | & + & 0 & - & | & +}}}{%
\begin{tikzpicture}[x=.48cm,y=.48cm,font=\scriptsize]
  \rooster{-5}{1}{5}{-3}{1}{3}
  \draw[grijs,dashed] (-2,-3) -- (-2,3.3) (2,-3) -- (2,3.3);
  \kromme{-5}{-2.18}{(\x-1)/((\x)^2-4)}
  \kromme{-1.83}{1.87}{(\x-1)/((\x)^2-4)}
  \kromme{2.06}{5}{(\x-1)/((\x)^2-4)}
  \vol{(1,0)}
\end{tikzpicture}}}

\begin{werkwijze}
\textbf{1.} Domein. \textbf{2.} Ontbind en vereenvoudig: perforaties. \textbf{3.} Nulwaarden. \textbf{4.} Verticale asymptoten. \textbf{5.} Horizontale of schuine asymptoot. \textbf{6.} Tekentabel met nulwaarden en polen. \textbf{7.} Een paar functiewaarden, en schetsen: de grafiek nadert de asymptoten aan de kant die de tekentabel aangeeft.
\end{werkwijze}

Een functie van de vorm $\frac{a}{x - p} + c$ is de grafiek van $\frac1x$, getransformeerd (Algemene functies): asymptoten $x = p$ en $y = c$. Elke functie $\frac{\alpha x + \beta}{\gamma x + \delta}$ (met $\gamma \neq 0$) is zo te schrijven, door te delen.

\begin{oefeningen}
\needspace{6\baselineskip}
\oef{h2-tekentabel}

\oef{h2-schets}

\needspace{6\baselineskip}
\oef{h2-homografisch}
\end{oefeningen}

%========================================================================================
\opdrachtenkop

\begin{oefeningen}
\needspace{14\baselineskip}
\oef{h2-o1}

\needspace{12\baselineskip}
\oef{h2-o10}

\needspace{16\baselineskip}
\oef{h2-o2}

\needspace{10\baselineskip}
\oef{h2-o12}

\needspace{6\baselineskip}
\oef{h2-o3}

\needspace{6\baselineskip}
\oef{h2-o11}

\needspace{8\baselineskip}
\oef{h2-o9}

\needspace{12\baselineskip}
\oef{h2-o14}

\oef{h2-o4}

\needspace{8\baselineskip}
\oef{h2-o19}

\needspace{12\baselineskip}
\oef{h2-o17}

\oef{h2-o5}

\needspace{6\baselineskip}
\oef{h2-o20}

\needspace{8\baselineskip}
\oef{h2-o21}

\needspace{6\baselineskip}
\oef{h2-o15}

\needspace{8\baselineskip}
\oef{h2-o22}

\needspace{9\baselineskip}
\oef{h2-o18}

\oef{h2-o6}

\needspace{12\baselineskip}
\oef{h2-o23}

\needspace{12\baselineskip}
\oef{h2-o16}

\oef{h2-o7}

\needspace{10\baselineskip}
\oef{h2-o25}

\needspace{10\baselineskip}
\oef{h2-o24}

\oef{h2-o27}

\needspace{8\baselineskip}
\oef{h2-o26}

\needspace{9\baselineskip}
\oef{h2-o8}
\end{oefeningen}

%----------------------------------------------------------------------------------------
% KOPPEL: verbanden met andere leerstof
%----------------------------------------------------------------------------------------
\needspace{20\baselineskip}
\opdrachtsectie{Koppel}

\lkr{Algemene functies, Afgeleiden en Integralen: kijk na wat al gezien is. Het pakket legt de volgorde niet vast.}

\begin{oefeningen}
\oef{h2-koppel-1}

\oef{h2-koppel-2}

\oef{h2-koppel-3}
\end{oefeningen}

%----------------------------------------------------------------------------------------
\opdrachtsectie{Test jezelf}

\begin{oefeningen}
\oef{h2-t1}

\oef{h2-t2}

\oef{h2-t5}

\oef{h2-t3}

\oef{h2-t4}

\oef{h2-t6}
\end{oefeningen}

\herhaalschema{%
  \fv{oef:h2-t1}{vb:h2-perforatie}{oef:h2-domein}%
  \fv{oef:h2-t2}{vb:h2-ha}{oef:h2-ha}%
  \fv{oef:h2-t5}{vb:h2-ha}{oef:h2-homografisch}%
  \fv{oef:h2-t3}{vb:h2-sa}{oef:h2-sa}%
  \fv{oef:h2-t4}{vb:h2-schets}{oef:h2-schets}%
  \fv{oef:h2-t6}{vb:h2-ha}{oef:h2-o7}}

\ifzelfstudie
\section*{Maak je eigen samenvatting}
Sluit dit boek en schrijf in je schrift: wat een rationale functie is, hoe je domein, nulwaarden en perforaties bepaalt, de regels voor verticale, horizontale en schuine asymptoten, waarom de deling de schuine asymptoot geeft, en de zeven stappen om de grafiek te schetsen. Vergelijk daarna met de kaders.
\fi
